\chapter{Charla Conferencia Jesús Garre Navarro}

La ponencia ofreció una visión amplia sobre la evolución del sector de la consultoría y sobre las competencias necesarias para desenvolverse en la llamada \emph{Imagination Age}. Supuso una oportunidad para conectar los contenidos trabajados en clase con ejemplos prácticos y permitió reflexionar sobre cómo orientar la carrera profesional en un entorno globalizado y en transformación constante.  

Al inicio de la intervención, mientras repasaba su trayectoria hasta la actualidad, el ponente subrayó la importancia de mantener aficiones fuera del ámbito laboral, utilizando su pasión por las motos como ejemplo de equilibrio personal. Este recordatorio resulta especialmente relevante en carreras exigentes, donde es fácil medirlo todo en términos de logros. En mi caso, la práctica continuada de la música a través del clarinete representa ese espacio de desconexión y disciplina que complementa la formación técnica. Esta idea enlaza con lo trabajado en clase sobre conciliación y calidad de vida laboral, subrayando que el bienestar personal constituye un factor clave para sostener el rendimiento en el tiempo.  

La ponencia puso también el foco en la dimensión global de la consultoría. Se destacó cómo gran parte del sector se articula en torno a la relación entre Estados Unidos e India, que constituye el eje principal de la internacionalización. Al mismo tiempo, se señaló que Europa mantiene un papel más rezagado, mientras que China avanza a gran velocidad y adquiere un protagonismo cada vez mayor en el mercado. El ponente comentó que había estado recientemente en China y que le había impresionado su dinamismo. Esta observación me resultó cercana, ya que durante el pasado verano también tuve la oportunidad de viajar al país y coincidí con esa percepción de un entorno vibrante y en constante transformación.  

El ponente expuso además la evolución de las distintas etapas recientes: de la \emph{Internet Age} a la \emph{Information Age}, la \emph{Digital Singularity} y la actual \emph{Imagination Age}. En este marco, la inteligencia artificial se presenta como una herramienta destinada a asumir tareas repetitivas, liberando a las personas para actividades creativas e innovadoras. Esta visión coincide con lo visto en clase sobre digitalización y futuro del trabajo, donde se remarcaba que la automatización redefine los roles y exige nuevas competencias.  

Entre las competencias señaladas como clave se encuentran la gestión del estrés, la resolución de problemas, el trabajo en entornos multiculturales, la orientación a resultados y la capacidad de aportar valor real al cliente. Estas cuestiones coinciden con lo trabajado en clase acerca de la importancia de la inteligencia emocional y la gestión de la diversidad.  

Por último, la ponencia invitó a reflexionar sobre las motivaciones que guían las decisiones profesionales: la estabilidad, el prestigio, el impacto social, la autonomía o el aprendizaje continuo. La referencia a los arquetipos de motivación laboral (operadores, artesanos, luchadores, donantes, exploradores y pioneros) coincidió con lo trabajado en clase acerca de la necesidad de alinear motivaciones individuales con los valores de la organización. Reconocer en qué perfil encaja cada persona constituye un ejercicio útil para comprender las decisiones de carrera y la manera en que se aporta valor en distintos contextos profesionales.  

En conjunto, la ponencia permitió relacionar la teoría estudiada en la asignatura con la experiencia práctica de un directivo del sector de la consultoría internacional. Más allá de presentar los retos actuales, aportó herramientas conceptuales para reflexionar sobre cómo orientar la carrera profesional en un escenario marcado por la globalización, la diversidad y la transformación tecnológica.  
